\documentclass[12pt, %
openright, 
oneside, %
%twoside, %TCC: Se seu texto tem mais de 100 páginas, descomente esta linha e comente a anterior
a4paper,    %
%english,   %
brazil]{facom-ufu-abntex2}

\usepackage{graphicx}
\graphicspath{{figuras/}{pictures/}{images/}{./}} % where to search for the images

\newcommand{\blue}[1]{\textcolor{blue}{#1}}
\newcommand{\red}[1]{\textcolor{red}{#1}}


\autor{Matheus José} %TCC
\data{2022}
\orientador{Renato Pimentel} %TCC
%\coorientador{Algum?} %TCC

% ---
% Informações de dados para CAPA e FOLHA DE ROSTO
% ---

\titulo{Dispensa de Disciplinas} %TCC

\hypersetup{pdfkeywords={palavra 1}{palavra 2}{palavra 4}{palavra 4}{palavra 5}} %TCC

\begin{document} 
\frenchspacing 

% ----------------------------------------------------------
% ELEMENTOS PRÉ-TEXTUAIS
% ----------------------------------------------------------
%\pretextual
\imprimircapa
\imprimirfolhaderosto


% ---
% Inserir folha de aprovação
% ---
%
% \includepdf{folhadeaprovacao_final.pdf} %TCC: depois de aprovado o trabalho, descomente esta linha e comente o próximo bloco para incluir scan da folha de aprovação.
%
\begin{folhadeaprovacao}

  \begin{center}
    {\ABNTEXchapterfont\large\imprimirautor}

    \vspace*{\fill}\vspace*{\fill}
    {\ABNTEXchapterfont\bfseries\Large\imprimirtitulo}
    \vspace*{\fill}
    
    \hspace{.45\textwidth}
    \begin{minipage}{.5\textwidth}
        \imprimirpreambulo
    \end{minipage}%
    \vspace*{\fill}
   \end{center}
    
   Trabalho em andamento. \imprimirlocal, 22 de junho de 2022: %TCC:

   \assinatura{\textbf{\imprimirorientador} \\ Renato Pimentel}  
   \assinatura{\textbf{Professor}}% \\ Convidado 1} %TCC:
   \assinatura{\textbf{Professor}}% \\ Convidado 2} %TCC:
   %\assinatura{\textbf{Professor} \\ Convidado 3}
   %\assinatura{\textbf{Professor} \\ Convidado 4}
      
   \begin{center}
    \vspace*{0.5cm}
    {\large\imprimirlocal}
    \par
    {\large\imprimirdata}
    \vspace*{1cm}
  \end{center}
  
\end{folhadeaprovacao}
% ---


%%As seções dedicatória, agradecimento e epígrafe não são obrigatórias.
%%Só as mantenha se achar pertinente.

% ---
% Dedicatória
% ---
%\begin{dedicatoria}
%   \vspace*{\fill}
%   \centering
%   \noindent
%   \textit{Dedico a \lipsum[10]}  %TCC:
%   \vspace*{\fill}
%\end{dedicatoria}
% ---

% ---
% Agradecimentos
% ---
%\begin{agradecimentos}
%Agradeço a \lipsum[30]. %TCC:
%\end{agradecimentos}
% ---

% ---
% Epígrafe
% ---
%\begin{epigrafe}
%    \vspace*{\fill}
%	\begin{flushright}
%		\textit{``Alguma citação que ache conveniente? \lipsum[10]''} %TCC:
%	\end{flushright}
%\end{epigrafe}
% ---



\begin{resumo} %TCC:
 Segundo a \citeonline[3.1-3.2]{NBR6028:2003}, o resumo deve ressaltar o
 objetivo, o método, os resultados e as conclusões do documento. A ordem e a extensão
 destes itens dependem do tipo de resumo (informativo ou indicativo) e do
 tratamento que cada item recebe no documento original. O resumo deve ser
 precedido da referência do documento, com exceção do resumo inserido no
 próprio documento. (\ldots) As palavras-chave devem figurar logo abaixo do
 resumo, antecedidas da expressão Palavras-chave:, separadas entre si por
 ponto e finalizadas também por ponto.

 \vspace{\onelineskip}
    
 \noindent
 \textbf{Palavras-chave}: Até, cinco, palavras-chave, separadas, por, vírgulas. %TCC:
\end{resumo}

% ---
% inserir lista de ilustrações
% ---
\pdfbookmark[0]{\listfigurename}{lof}
\listoffigures*
\cleardoublepage
% ---

% ---
% inserir lista de tabelas
% ---
\pdfbookmark[0]{\listtablename}{lot}
\listoftables*
\cleardoublepage
% ---



% ---
% inserir lista de abreviaturas e siglas
% ---
\begin{siglas} %TCC:
  \item[Fig.] Area of the $i^{th}$ component
  \item[456] Isto é um número
  \item[123] Isto é outro número
  \item[Zézão] este é o meu nome
\end{siglas}
% ---

%% ---
%% inserir lista de símbolos, se for adequado ao trabalho. %TCC:
%% ---
%\begin{simbolos}
%  \item[$ \Gamma $] Letra grega Gama
%  \item[$ \Lambda $] Lambda
%  \item[$ \zeta $] Letra grega minúscula zeta
%  \item[$ \in $] Pertence
%\end{simbolos}
%% ---

% ---
% inserir o sumario
% ---
\pdfbookmark[0]{\contentsname}{toc}
\tableofcontents*
\cleardoublepage
% ---





% ----------------------------------------------------------
% ELEMENTOS TEXTUAIS
% ----------------------------------------------------------
\textual


% ----------------------------------------------------------
% Introdução
% ----------------------------------------------------------

\chapter[Introdução]{Introdução}
%TCC:
Com a crescente entrada de pessoas à universidade, as coordenações de cursos têm de lidar com as mais diversas situações da vida escolar dos alunos. Pode-se citar, por exemplo, jovens que iniciam em determinado curso, porém não se sentem tão confortáveis e acabam migrando para outro curso ou até mesmo para outra instituição de ensino, e que porventura abrem solicitação para a coordenação para avaliar equivalência de disciplinas.

Essas solicitações que alunos encaminham às coordenações acabam gerando bastante gargalo de processos, justamente porque existem inúmeros cenários distintos na análise de equivalência.

Hoje, na Universidade Federal de Uberlândia, quando um aluno requisita equivalência de disciplinas para a coordenação de curso é passado para um professor que tenha conhecimento da ementa daquela disciplina. Esse compara a ementa que o aluno trouxe – do curso em que ele concluiu aquela matéria – com o curso que ele está. Essa parte é a mais minuciosa, pois ela requer uma análise crítica de cada componente da ementa. E então é dado um parecer da análise. Se aprovado, é emitido uma certidão comprobatória da equivalência. Caso não seja aceito, o aluno então poderá recorrer em um prazo estipulado. 

Essa fase de análise de ementa é a que mais gasta esforço humano, e por essa razão surgiu a implementação de um sistema que seja capaz de analisar se, entre duas ementas de disciplinas, essas sendo da mesma universidade ou não, possuem equivalência.

Para que o sistema funcione, será necessário uma página na internet para cadastro, listagem, deleção e atualização das disciplinas – com suas respectivas ementas –, isso tudo separado por curso e universidade.

Com esses dados, então será possível analisar, com base na ementa, se entre duas disciplinas, essas são equivalentes.

Além disso, outra forma de adicionar dados ao sistema será a partir da leitura de um arquivo PDF e também de imagem. Essa funcionalidade será para diminuir o tempo de inserção de dados manuais.

Dessa forma a análise da equivalência de disciplinas se tornará algo menos custoso operacionalmente, pois a pessoa que está analisando receberá uma resposta sistêmica que servirá como uma pré-aprovação da equivalência. 
\cite{ibge1993}

\chapter{Revisão Bibliográfica}
%TCC:
Um ou mais capítulos (por exemplo, se há duas linhas de trabalhos relacionados).


\chapter{Desenvolvimento}
%TCC:
Um ou mais capítulos (por exemplo um para testes)


\begin{figure}[!ht]
    \centering
	\includegraphics[width=0.55\linewidth]{imagemExemplo.pdf}
	\caption[Isso é o que aparece no sumário]{Imagem de exemplo.}
	\label{fig:graficosVariandoTamanhoRede}
\end{figure}


%TCC:
%TCC:
%TCC:
%TCC:

% ---
% Conclusão
% ---
\chapter[Conclusão]{Conclusão}
%TCC:
E daí?


\chapter[Cronograma]{Cronograma}
  \begin{table}[htbp]
    \centering
    \resizebox{\textwidth}{!}{
    \begin{tabular}{c|c|c|c|c|c|c|c|c|c|c|}
    \hline
    ~ & Junho & Julho & Agosto & Setembro & Outubro & Novembro & Dezembro & Dezembro & Janeiro & Fevereiro\\
    \hline
    \hline
    Levantamento de requisitos & X & X & ~ & ~ & ~ & ~ & ~ & ~ & ~ & ~ \\
    Definição de tecnologias  & X & ~ & ~ & ~ & ~ & ~ & ~ & ~ & ~ & ~ \\
    Acompanhamento de disciplina PG1 & X & X & X & ~ & ~ & ~ & ~ & ~ & ~ & ~ \\
    Desenvolvimento banco de dados & ~ & X & ~ & ~ & ~ & ~ & ~ & ~ & ~ & ~ \\
    Backend:Desenvolvimento cadastro de disciplinas & ~ & X & X & X & ~ & ~ & ~ & ~ & ~ & ~ \\
    Backend:Avaliar equivalência entre duas disciplinas & ~ & X & X & X & ~ & ~ & ~ & ~ & ~ & ~ \\
    Backend:Cadastro de equivalências  & ~ & ~ & X & X & ~ & ~ & ~ & ~ & ~ & ~ \\
    Backend:Desenvolvimento de cadastro de disciplinas por leitura de PDF/imagem & ~ & ~ & ~ & X & X & ~ & ~ & ~ & ~ & ~ \\
    Frontend:Desenvolvimento da tela de cadastro & ~ & ~ & ~ & X & X & X & ~ & ~ & ~ & ~ \\
    Frontend:Desenvolvimento da tela de verificacao se dá equivalencia ou não & ~ & ~ & ~ & X & X & X & ~ & ~ & ~ & ~ \\
    Frontend:Desenvolvimento da tela com retorno se deu equivalencia ou não & ~ & ~ & ~ & X & X & X & ~ & ~ & ~ & ~ \\
    Testes unitários / Funcionais & ~ & ~ & X & X & X & ~ & ~ & ~ & ~ & ~ \\
    Escrita texto & ~ & ~ & ~ & X & X & X & X & ~ & ~ & ~ \\
    Revisão texto & ~ & ~ & ~ & ~ & ~ & X & X & X & X & ~ \\
    Defesa & ~ & ~ & ~ & ~ & ~ & ~ & ~ & ~ & ~ & X \\

    \end{tabular}
    }
    \caption{Cronograma}
    \label{tab:my_label}
  \end{table}


% ----------------------------------------------------------
% ELEMENTOS PÓS-TEXTUAIS
% ----------------------------------------------------------
\postextual


% ----------------------------------------------------------
% Referências bibliográficas
% ----------------------------------------------------------
\bibliography{abntex2-modelo-references}


%% ----------------------------------------------------------
%% Apêndices TCC: só mantenha se for pertinente.
%% ----------------------------------------------------------

% ---
% Inicia os apêndices
% ---
\begin{apendicesenv}

% Imprime uma página indicando o início dos apêndices
\partapendices

% ----------------------------------------------------------
\chapter{Quisque libero justo}
% ----------------------------------------------------------

\lipsum[50]

% ----------------------------------------------------------
\chapter{Coisas que fiz e que achei interessante mas não tanto para entrar no corpo do texto}
% ----------------------------------------------------------
\lipsum[55-57]

\end{apendicesenv}
% ---


% ----------------------------------------------------------
% Anexos %TCC: so mantenha se pertinente.
% ----------------------------------------------------------

% ---
% Inicia os anexos
% ---
\begin{anexosenv}

% Imprime uma página indicando o início dos anexos
\partanexos

% ---
\chapter{Eu sempre quis aprender latim}
% ---
\lipsum[30]

% ---
\chapter{Coisas que eu não fiz mas que achei interessante o suficiente para colocar aqui}
% ---

\lipsum[31]

% ---
\chapter{Fusce facilisis lacinia dui}
% ---

\lipsum[32]

\end{anexosenv}

%---------------------------------------------------------------------
% INDICE REMISSIVO
%---------------------------------------------------------------------

\printindex



\end{document}